\let\textcircled=\pgftextcircled
\chapter{Travail réalisé, limites et perspectives}
\label{chap:results}

\initial{C}\textit{e chapitre distingue ce qui a été livré en deux mois de ce qui relève du protocole encore à exécuter sur les extraits, et énonce les limites honnêtes du stage.}

\newpage 

\section{Livrables du mois 1}

Le mois 1 (24 juillet -- fin août 2026) a produit le socle scientifique et administratif, avant les modèles :

\begin{itemize}[label=\ding{42}, font=\large \color{darkorange}]
\item Cartographie de huit bases (NACC, ADNI, Framingham, OASIS-3, EPAD, AMYPAD, AIBL, UK Biobank) : modalités, volumes, DUA, délais, IRB.
\item État de l'art sur la fusion multimodale et les modalités manquantes (présentations de travail).
\item Revue de littérature ciblée : CARD/NIH, Xue--Kolachalama 2024, MOIRA 2025, fuites, PROBAST, calibration/\ac{DCA}, \ac{SHAP}.
\item Plan d'analyse rédigé sur le canevas Fongang (anglais pour Dr Fongang, français pour le rapport interne), avec trois aims alignés et calendrier recalé sur le 24 juillet -- 24 septembre.
\end{itemize}

Ces pièces répondent à la consigne initiale : << commencer par trouver les jeux de données ; proposer un plan of analytics >>.

\section{Calendrier réel (deux mois)}

\begin{center}
\begin{tabular}{p{0.55\linewidth}p{0.35\linewidth}}
\hline
\textbf{Tâche} & \textbf{Période} \\
\hline
Cartographie et voies d'accès & 24 jul.\ -- 7 août \\
État de l'art fusion / manques & 8--20 août \\
Revue CARD / MOIRA / Kolachalama & 21--25 août \\
Brouillon du plan d'analyse & 25--26 août \\
Plan final après commentaires & 1$^{\mathrm{er}}$ septembre \\
Inventaire des extraits, lock phénotype & 26 août -- 3 septembre \\
Table d'analyse et splits participant & 3--8 septembre \\
Modèles 1--2 puis 3--4, calibration, DCA & 8--17 septembre \\
Validation externe, SHAP, équité, paquet minimal & 17--21 septembre \\
Rapport et restitution & 22--24 septembre \\
\hline
\end{tabular}
\end{center}

Le goulot n'est plus une file NACC/LONI/BioLINCC \emph{si} les extraits mentors sont effectivement disponibles. Sinon, la validation externe réaliste se réduit à NACC $\leftrightarrow$ ADNI, et Framingham est déclaré comme limite, pas comme échec.

\section{Ce que ce rapport ne contient pas encore}

À la date de rédaction, les tableaux démographiques (N MCI, taux de conversion à 2 et 5 ans, couverture IRM/TEP) restent vides : ils se remplissent \emph{après} inventaire des fichiers, pas par invention. Les AUC, C-index, courbes de calibration et \ac{DCA} empiriques seront insérés dès que les modèles 1--4 auront tourné sous DUA. Ce n'est pas un manquement de protocole : c'est la discipline même du plan d'analyse (pas de résultats avant phénotype locké).

\section{Limites}

\begin{itemize}[label=\ding{42}, font=\large \color{darkorange}]
\item Durée de deux mois : pas de transformer type Kolachalama, pas de CNN 3D comme résultat principal.
\item NACC/ADNI sont des échantillons cliniques ; Framingham est peu divers historiquement. Les données africaines, prévues comme test de transportabilité ultérieur, sont hors scope.
\item L'équité ne peut être rapportée que si les variables et les effectifs de cellules existent.
\item Le DUA est nominatif : le stagiaire doit être ajouté à l'accord du laboratoire le cas échéant.
\end{itemize}

\section{Perspectives}

Une fois les extraits inventoriés : figer les codes MCI/démence avec les encadrants ; exécuter les modèles 1--4 ; publier le paquet minimal comme résultat principal, même si l'AUC interne est moins spectaculaire qu'un papier ADNI-only. À plus long terme : test de transportabilité sur données africaines harmonisées, dans le cadre AI-BOND.
